% Project-specific PDF typography. Loaded after Verso's generated preamble.
\usepackage{luatexko}
\usepackage{unicode-math}
\usepackage{etoolbox}
\usepackage[a4paper,margin=24mm,headheight=16pt]{geometry}

\directlua{
  luaotfload.add_fallback("bookmath", {
    "[fonts/NotoSansMath-Regular.otf]:mode=node;",
    "[fonts/Pretendard-Regular.ttf]:mode=node;"
  })
}

\setmainfont{Pretendard-Regular.ttf}[
  Path=fonts/,
  BoldFont=Pretendard-Bold.ttf,
  ItalicFont=Pretendard-Regular.ttf,
  ItalicFeatures={FakeSlant=0.15},
  BoldItalicFont=Pretendard-Bold.ttf,
  BoldItalicFeatures={FakeSlant=0.15},
  RawFeature={fallback=bookmath}
]
\setsansfont{Pretendard-Regular.ttf}[
  Path=fonts/,
  BoldFont=Pretendard-Bold.ttf,
  ItalicFont=Pretendard-Regular.ttf,
  ItalicFeatures={FakeSlant=0.15},
  BoldItalicFont=Pretendard-Bold.ttf,
  BoldItalicFeatures={FakeSlant=0.15},
  RawFeature={fallback=bookmath}
]
\setmainhangulfont{Pretendard-Regular.ttf}[
  Path=fonts/,
  BoldFont=Pretendard-Bold.ttf,
  ItalicFont=Pretendard-Regular.ttf,
  ItalicFeatures={FakeSlant=0.15},
  BoldItalicFont=Pretendard-Bold.ttf,
  BoldItalicFeatures={FakeSlant=0.15}
]
\setsanshangulfont{Pretendard-Regular.ttf}[
  Path=fonts/,
  BoldFont=Pretendard-Bold.ttf,
  ItalicFont=Pretendard-Regular.ttf,
  ItalicFeatures={FakeSlant=0.15},
  BoldItalicFont=Pretendard-Bold.ttf,
  BoldItalicFeatures={FakeSlant=0.15}
]
\setmonofont{CascadiaCode-Regular.ttf}[
  Path=fonts/,
  BoldFont=CascadiaCode-Bold.ttf,
  ItalicFont=CascadiaCodeItalic-Regular.ttf,
  BoldItalicFont=CascadiaCodeItalic-Bold.ttf,
  Ligatures=NoCommon,
  RawFeature={-calt,fallback=bookmath}
]
\setmonohangulfont{Pretendard-Regular.ttf}[
  Path=fonts/,
  BoldFont=Pretendard-Bold.ttf,
  ItalicFont=Pretendard-Regular.ttf,
  ItalicFeatures={FakeSlant=0.15},
  BoldItalicFont=Pretendard-Bold.ttf,
  BoldItalicFeatures={FakeSlant=0.15}
]
\setmathfont{TeX Gyre Pagella Math}
\setmathfont{NotoSansMath-Regular.otf}[
  Path=fonts/,
  Scale=1.10,
  range={"00AC,"00B1,"00D7,"00F7,"2190-"22FF,
    "27C0-"27FF,"2900-"2AFF}
]

% Override symbols even when Cascadia contains them. A missing-glyph
% fallback alone leaves arrows such as U+2194 in the code font.
\newfontfamily\bookSymbolFont{NotoSansMath-Regular.otf}[
  Path=fonts/,
  Scale=1.10,
  BoldFont=NotoSansMath-Regular.otf,
  ItalicFont=NotoSansMath-Regular.otf,
  BoldItalicFont=NotoSansMath-Regular.otf
]
\newcommand{\BookMathSymbol}[1]{{\bookSymbolFont\char"#1\relax}}
\newcommand{\BookSymbolCodes}{%
  \catcode"00AC=\active\relax
  \catcode"00B1=\active\relax
  \catcode"00D7=\active\relax
  \catcode"00F7=\active\relax
  \catcode"2115=\active\relax
  \catcode"2124=\active\relax
  \catcode"2190=\active\relax
  \catcode"2192=\active\relax
  \catcode"2194=\active\relax
  \catcode"21A6=\active\relax
  \catcode"21D2=\active\relax
  \catcode"21D4=\active\relax
  \catcode"2200=\active\relax
  \catcode"2203=\active\relax
  \catcode"2205=\active\relax
  \catcode"2208=\active\relax
  \catcode"2209=\active\relax
  \catcode"2218=\active\relax
  \catcode"2223=\active\relax
  \catcode"2224=\active\relax
  \catcode"2227=\active\relax
  \catcode"2228=\active\relax
  \catcode"2229=\active\relax
  \catcode"222A=\active\relax
  \catcode"2234=\active\relax
  \catcode"2260=\active\relax
  \catcode"2261=\active\relax
  \catcode"2264=\active\relax
  \catcode"2265=\active\relax
  \catcode"2282=\active\relax
  \catcode"2283=\active\relax
  \catcode"2286=\active\relax
  \catcode"2287=\active\relax
  \catcode"22A2=\active\relax
  \catcode"27E8=\active\relax
  \catcode"27E9=\active\relax
}
\begingroup
\BookSymbolCodes
\protected\gdef¬{\BookMathSymbol{00AC}}
\protected\gdef±{\BookMathSymbol{00B1}}
\protected\gdef×{\BookMathSymbol{00D7}}
\protected\gdef÷{\BookMathSymbol{00F7}}
\protected\gdefℕ{\BookMathSymbol{2115}}
\protected\gdefℤ{\BookMathSymbol{2124}}
\protected\gdef←{\BookMathSymbol{2190}}
\protected\gdef→{\BookMathSymbol{2192}}
\protected\gdef↔{\BookMathSymbol{2194}}
\protected\gdef↦{\BookMathSymbol{21A6}}
\protected\gdef⇒{\BookMathSymbol{21D2}}
\protected\gdef⇔{\BookMathSymbol{21D4}}
\protected\gdef∀{\BookMathSymbol{2200}}
\protected\gdef∃{\BookMathSymbol{2203}}
\protected\gdef∅{\BookMathSymbol{2205}}
\protected\gdef∈{\BookMathSymbol{2208}}
\protected\gdef∉{\BookMathSymbol{2209}}
\protected\gdef∘{\BookMathSymbol{2218}}
\protected\gdef∣{\BookMathSymbol{2223}}
\protected\gdef∤{\BookMathSymbol{2224}}
\protected\gdef∧{\BookMathSymbol{2227}}
\protected\gdef∨{\BookMathSymbol{2228}}
\protected\gdef∩{\BookMathSymbol{2229}}
\protected\gdef∪{\BookMathSymbol{222A}}
\protected\gdef∴{\BookMathSymbol{2234}}
\protected\gdef≠{\BookMathSymbol{2260}}
\protected\gdef≡{\BookMathSymbol{2261}}
\protected\gdef≤{\BookMathSymbol{2264}}
\protected\gdef≥{\BookMathSymbol{2265}}
\protected\gdef⊂{\BookMathSymbol{2282}}
\protected\gdef⊃{\BookMathSymbol{2283}}
\protected\gdef⊆{\BookMathSymbol{2286}}
\protected\gdef⊇{\BookMathSymbol{2287}}
\protected\gdef⊢{\BookMathSymbol{22A2}}
\protected\gdef⟨{\BookMathSymbol{27E8}}
\protected\gdef⟩{\BookMathSymbol{27E9}}
\endgroup

\linespread{1.28}
\setlength{\parindent}{0pt}
\setlength{\parskip}{0.45em}
\setcounter{tocdepth}{1}
\setcounter{secnumdepth}{3}
\renewcommand{\contentsname}{목차}
\renewcommand{\chaptername}{장}
\renewcommand{\figurename}{그림}
\renewcommand{\tablename}{표}

% Keep the four-column symbol reference within the text area.
\newlength{\BookSymbolTableWidth}
\newenvironment{BookSymbolTable}{%
  \par\begingroup
  \linespread{1.0}%
  \fontsize{9.5}{11.5}\selectfont
  \setlength{\parskip}{0pt}%
  \setlength{\tabcolsep}{4pt}%
  \setlength{\BookSymbolTableWidth}{\dimexpr\linewidth-6\tabcolsep\relax}%
  \renewcommand{\arraystretch}{1.08}%
  \begin{tabular}{@{}%
    >{\raggedright\arraybackslash}p{0.22\BookSymbolTableWidth}%
    >{\raggedright\arraybackslash}p{0.09\BookSymbolTableWidth}%
    >{\raggedright\arraybackslash}p{0.36\BookSymbolTableWidth}%
    >{\raggedright\arraybackslash}p{0.33\BookSymbolTableWidth}@{}}%
}{\end{tabular}\endgroup\par}

\definecolor{codeBorder}{HTML}{D6DCE3}
\DefineVerbatimEnvironment{LeanVerbatim}{Verbatim}{
  commandchars=\\\{\},
  codes*={\BookSymbolCodes},
  fontsize=\small,
  baselinestretch=1.35,
  breaklines=true,
  breakanywhere=true,
  samepage=false,
  frame=single,
  rulecolor=\color{codeBorder},
  framesep=2mm,
  vspace=0pt,
  xleftmargin=1mm,
  xrightmargin=1mm
}
\RecustomVerbatimCommand{\LeanVerb}{Verb}{
  commandchars=\\\{\},codes*={\BookSymbolCodes},
  fontsize=\small,breaklines=true,breakanywhere=true
}
\BeforeBeginEnvironment{LeanVerbatim}{%
  \par\begingroup
  \setlength{\parskip}{0pt}%
  \setlength{\topsep}{0pt}%
  \setlength{\partopsep}{0pt}%
}
\AfterEndEnvironment{LeanVerbatim}{\endgroup\par}

% Memoir uses slanted text in its running headings.
\makeatletter
\DeclareFontShape{TU}{\rmdefault}{m}{sl}{<->ssub * \rmdefault/m/it}{}
\makeatother

\hypersetup{
  unicode=true,
  colorlinks=true,
  linkcolor={blue!45!black},
  urlcolor={blue!45!black},
  pdftitle={Lean으로 배우는 기초 수학},
  pdfsubject={한국어 대수학 교재와 Lean 예제}
}
\pretitle{\begin{center}\sffamily\bfseries\fontsize{30}{40}\selectfont}
\posttitle{\par\end{center}}

% Keep headings with the paragraph that follows them.
\clubpenalty=10000
\widowpenalty=10000
\displaywidowpenalty=10000
\emergencystretch=2em
